\subsection{Independent core, tail and upper maps}
\label{rw:source:construction}

For each fixed field choose the reflection convention of
\eqref{rw:geometry:reflection} so that $\ell\ge0$, and retain this
direction throughout the count. The analytic constructions use signed
independent forcing parameters before this fixed-field restriction.
For a central height $s$ put
\begin{equation}\label{rw:source:scales}
 h=s-1-\log s,\quad L=-\log s,\quad
 \vartheta=h^{-1/2},\quad\widehat\ell=\ell/s,\qquad
 \beta=(Bh,mh,a\sqrt h,c\sqrt h).
\end{equation}
For the construction and coefficient division, first regard
$\vartheta,s,\beta,\widehat\ell$ as independent, with
\begin{equation}\label{rw:source:independent-box}
 0<\vartheta<\vartheta_0,\qquad
 0<s\le e^{-1/\vartheta^2},\qquad
 |\widehat\ell|\le3,\qquad |\beta_i|\le b_0.
\end{equation}
Here $b_0>0$ is fixed sufficiently small and $\vartheta_0>0$
is fixed afterwards. A selected jet is a word in
$E_\vartheta=\vartheta\partial_\vartheta$, $E_s=s\partial_s$,
$\partial_{\beta_i}$, $\partial_{\widehat\ell}$, and ordinary independent
section derivatives. Every estimate refers to a prescribed finite
number of such derivatives. The physical relations in
\eqref{rw:source:scales} are imposed after
these independent estimates. Since $\widehat\ell=\ell/s$, their
weighted forcing jets give no bound for unweighted $\partial_\ell$.

On the outgoing core use $X=\vartheta x$, $0\le X\le1$, and set
\begin{equation}\label{rw:source:core-data}
\begin{split}
 F(X)&=1+\beta_BX^2+\beta_aX+\vartheta^2\beta_m,\\
 D(X,z)&=F(X)-(1+2\vartheta^2\beta_m)z+\vartheta^2\beta_mz^2,\\
 J(X)&=\exp\int_0^X\frac{t/\vartheta^2-\beta_c}{F(t)}\,dt,
 \qquad I(X)=\int_0^X\frac{dt}{\vartheta F(t)J(t)}.
\end{split}
\end{equation}
The exact original graph equation is
\begin{equation}\label{rw:source:core-equation}
 z_X=\frac{s\widehat\ell+(X/\vartheta-\vartheta\beta_c-2s\widehat\ell)z
                      +(\vartheta\beta_c+s\widehat\ell)z^2}
                 {\vartheta D(X,z)},\qquad z(0)=s.
\end{equation}
The incoming core is the same construction after replacing
$(\beta_a,\beta_c,\widehat\ell)$ by $(-\beta_a,-\beta_c,-\widehat\ell)$.
Denote their endpoint quantities by subscripts $+$ and $-$.

\begin{proposition}[Source transition maps]\label{rw:source-maps}
After the preceding fixed choices the following statements hold on
common neighborhoods of their real domains.
\begin{enumerate}[label=(\roman*),leftmargin=*]
\item The signed normalized solution $U=z/(sJ)$ extends through
$z=0$ and satisfies
\begin{equation}\label{rw:source:core-expansion}
 \begin{gathered}
 U(X)=1+\widehat\ell I(X)+R(X),\qquad
 |R|_N\le C_N e^{-1/(3\vartheta^2)},\\
 |U|<10,
 \quad |z|<1/8,\quad D>27/32.
 \end{gathered}
\end{equation}
The selected jets of $I$ are bounded, $4/9<I(1)<2$, and
$U_{\widehat\ell}(1)>1/4$. There is a unique simple threshold
$\widehat\ell_c\in(-3,0)$; the outgoing core stays positive through
$X=1$ exactly when $\widehat\ell>\widehat\ell_c$. At equality its first zero
is at $X=1$.
\item For an independent entry $v\in[7/16,9/16]$, the substitution
$z=\exp[-(1-r)v/\vartheta^2]$ gives the tail equation
\begin{equation}\label{rw:source:tail-equation}
 X_r=v\frac{1-z+\beta_BX^2+\beta_aX+
                    \vartheta^2\beta_m(1-z)^2}
 {X+\vartheta^2\beta_c(z-1)+\vartheta s\widehat\ell(1-z)^2/z},
 \qquad X(0)=1.
\end{equation}
It reaches $r=1$, hence $z=1$, in finite positive original time,
with $3/4<X<2$ and denominator $>3/4$.
Its action map $T(v)=X(1)^2/2$ has bounded selected jets,
$63/64<T_v<65/64$, and is uniformly close to $1/2+v$.
Its image contains the fixed interval $[31/32,33/32]$.
The inverse on that interval has the same jet bounds. Moreover
\begin{equation}\label{rw:source:tail-forcing}
 T_{\widehat\ell}-T_0=\widehat\ell H_T,\qquad
 |H_T|_N\le C_Ne^{-1/(3\vartheta^2)},
\end{equation}
and the corresponding inverse difference has the same divided bound.
\item With $Y=\vartheta^2y$, $\lambda=\vartheta^2$ and
$\alpha=s\widehat\ell/\vartheta$, the upper field is
\begin{equation}\label{rw:source:upper-field}
 X_\xi=Y(\beta_mY-1)+\lambda(\beta_BX^2+\beta_aX),\qquad
 Y_\xi=(Y+\lambda)(X+\beta_cY)+\alpha Y^2.
\end{equation}
For sufficiently small $|\alpha|+\sum|\beta_i|$, each action
$e=X^2/2\in[3/4,5/4]$ at $Y=0$, $X>0$, has a first upper
passage to $Y=0$, $X<0$. It remains in $|X|<2$, $0\le Y<3/2$.
Its action map $\mathcal U_\alpha(e)$ satisfies
$|\mathcal U_\alpha(e)-e|<1/128$, with bounded finite jets in
$e,\alpha,\beta,E_\lambda$.
\end{enumerate}
\end{proposition}
\begin{proof}
Appendix~\ref{rw:app:source-estimates} proves the independent-entry
jets and real passage conditions using the reciprocal Gaussian
$J^{-1}$ in the core, the exponentially small tail forcing $s/z$,
and the integrable upper carrier $\lambda/(Y+\lambda)$.
All three estimates precede substitution of logarithmic core values
into the tail entries.
\end{proof}

Now impose the physical relation \eqref{rw:source:scales}. Fix
$\delta_0=1/8192$ and retain the positive-core set
\begin{equation}\label{rw:source:positive-core}
 U_+\ge s^{2\delta_0},\qquad U_-\ge s^{2\delta_0}.
\end{equation}
The actual tail entries are
\[
 v_\pm=-h^{-1}\log(sJ_\pm(1)U_\pm).
\]
Since $h^{-1}\log J_\pm(1)=1/2+O(b_0)$,
$L/h=1+O(1/h)$, and $|U_\pm|<10$, they lie in a fixed small
neighborhood of $1/2$, strictly inside $[7/16,9/16]$.
Choose the coefficient box smaller once more and then $h$ large so
that their actions $e_\pm=T_{\pm,\widehat\ell}(v_\pm)$ satisfy
$|e_\pm-1|<1/1024$. Both $e_-$ and
$\mathcal U_\alpha(e_+)$ then belong to the common tail-inverse
interval. On this domain the exact shooting function is
\begin{equation}\label{rw:source:shooting}
 F_{\widehat\ell}=\mathcal U_\alpha(e_+)-e_-.
\end{equation}
Its zero joins the two independent lower arcs by the actual upper arc.
Action is injective on the negative $Y=0$ section, so this is exactly
the closing condition for these five pieces. Near a closing orbit the
backward incoming section map is an analytic diffeomorphism; its
nonzero derivative shows that \eqref{rw:source:shooting} has the
same zero multiplicity as the Poincare displacement. This argument
also applies at retained finite endpoints of
\eqref{rw:source:positive-core}, because all tube inequalities there
are strict. Existence of a positive forward central return throughout
this set is proved below in the fixed direction $\ell\ge0$.

\subsection{Center division and the physical derivative}

At $\widehat\ell=0$ write $L_\pm=T_{\pm,0}$,
$v_{\pm,0}=v_\pm|_{\widehat\ell=0}$, and define, on a fixed entry
interval about $1/2$,
\begin{equation}\label{rw:source:unforced-map}
 \mathcal M=L_-^{-1}\circ\mathcal U_0\circ L_+,
 \qquad G_0=\mathcal M(v_{+,0})-v_{-,0}.
\end{equation}
The common tail-inverse domain makes this composition well defined.
We now divide $G_0$ by the two center generators. The divided
coefficients will give a positive amplitude and an increasing energy
coordinate for the forced closing equation.

\begin{lemma}[Divided unforced source function]
\label{rw:source:center-division}
On a fixed smaller coefficient box,
\begin{equation}\label{rw:source:unforced-division}
 G_0=\tau_0\sqrt h\,A+\sigma_0h^{3/2}C,
 \qquad \tau_0=a+c,\quad\sigma_0=c(B+m).
\end{equation}
After $s=s(\vartheta)$ is imposed, $A,C$ have bounded selected
Euler and scaled-coefficient jets of every prescribed finite order.
Their limiting functions as $\vartheta\to0$ are analytic in $\beta$,
and their errors in these jets are
$O(r_s)$, where $r_s=\vartheta^2(1+|\log\vartheta|)$.
At $\beta=0$ their limits are
\begin{equation}\label{rw:source:moments}
 A_*={4\sqrt2}/{3},\qquad C_*={8\sqrt2}/{15},\qquad
 \mathcal M_{\rm pr}(v,0)=v.
\end{equation}
\end{lemma}
\begin{proof}
Use the scaled coefficient coordinates
$T=\beta_a+\beta_c$, $u=\beta_a$, $t=\beta_B+\beta_m$ and
$\beta_B$. The cube
$(T,u,t)\mapsto(q_1T,q_2u,q_3t)$, $0\le q_i\le1$, remains
inside the larger independent coefficient box after a fixed initial
shrinking. Indeed
$\beta_m'=q_3\beta_m-(1-q_3)\beta_B$ and
$\beta_c'=q_1\beta_c+(q_1-q_2)\beta_a$.
At $\widehat\ell=0$ both normalized cores are $1+R>1/2$, so the entire
cube has all the strict real passage margins.

The unforced shooting function vanishes on $T=u=0$ by reversibility:
reflection interchanges the lower arcs and the two halves of the
upper arc. It also vanishes on $T=t=0$. This is the second real
center of Proposition~\ref{rw:real-centers}, applied to the scaled
coefficients $(\beta_B,\beta_m,\beta_a,\beta_c)$ and scale
$\lambda=\vartheta^2$. Its polynomial $A_C$ is positive on
$[-2,2]$ in the small source box; hence the actual outgoing $Y=0$
point lies in its period annulus. Uniqueness identifies its periodic
arc with the constructed passages. Thus both center identities hold
on the whole contraction cube used for division.

If $\Psi(T,u,t,\beta_B)=F_0$, successive fundamental-theorem
integrals give
\begin{equation}\label{rw:source:Hadamard-rule}
\begin{gathered}
 \Psi=TM+utN,\quad
 M=\int_0^1\Psi_T(qT,u,t,\beta_B)\,dq,\\
 N=\int_0^1\!\int_0^1\Psi_{ut}(0,qu,rt,\beta_B)\,dq\,dr.
\end{gathered}
\end{equation}
Since $ut=Tt-\beta_ct$, the coefficients of
$T$ and $\beta_ct$ are $A_0=M+tN$ and $C_0=-N$.
The mean-value formula for $L_-^{-1}$ gives
$G_0=R_0F_0$, with
\[
 R_0=\int_0^1(L_-^{-1})'(e_{-,0}+qF_0)\,dq>0.
\]
Thus $A=R_0A_0$ and $C=R_0C_0$ give
\eqref{rw:source:unforced-division}. These specified integrals also
fix the divided coefficients on a compact overlap used later.

To justify physical substitution, note that
$E_\vartheta s=2hs/(1-s)$. Each fixed number of derivatives adds
only polynomial factors of $h$ to an exponentially small core term.
The exact unforced entries are
\[
 v_{\pm,0}=1+\vartheta^2(1-s)
 -\int_0^1\frac{t}{F_\pm(t)}\,dt
 +\vartheta^2\beta_{c,\pm}\int_0^1\frac{dt}{F_\pm(t)}
 -\vartheta^2\log(1+R_\pm).
\]
They consequently have the stated bounded jets. Two extra coefficient
orders before the Hadamard integrations supply any desired final order.

The limiting lower endpoints $\xi_+>0$, $\xi_-<0$ solve
\[
 S_\beta(\xi_\pm)=1,\qquad
 S_\beta(X)=\int_0^X\frac{t}{1+\beta_at+\beta_Bt^2}\,dt.
\]
The limiting upper system is $X'=-1+\beta_mY$,
$Y'=X+\beta_cY$. Appendix~\ref{rw:app:source-principal}
compares the exact and limiting integral equations in the reserved
jets, giving error $O(r_s)$ on the whole coefficient cube. The same
Hadamard formulas and common tail inverse give that error for $A,C$.

For the constants take $\beta_c=0$, when the upper first integral is
$X^2/2+Y-\beta_mY^2/2$. At input action $e$, put
$q_e=e-X^2/2$ and $a_e=\sqrt{2e}$. Along this arc
$Y=q_e+\tfrac12\beta_mq_e^2+O(\beta_m^2)$, and variation in
$\beta_c$ gives $\partial_{\beta_c}\mathcal U_{\rm pr}
=\int_{-a_e}^{a_e}Y\,dX$. Hence
\[
 \int_{-a_e}^{a_e}q_e\,dX=\frac{4\sqrt2}{3}e^{3/2},\qquad
 \frac12\int_{-a_e}^{a_e}q_e^2\,dX
                         =\frac{8\sqrt2}{15}e^{5/2}.
\]
The lower endpoints are independent of $\beta_m,\beta_c$ in this
limit. The same Hadamard rule, in particular $C_0=-N$, gives the
positive constants in \eqref{rw:source:moments}. At $\beta=0$ the
tail map is $v\mapsto1/2+v$ and $R_0=1$ in the limit.
\end{proof}

On the unforced coefficients, after physical substitution, define
\begin{equation}\label{rw:source:physical-E}
 \mathfrak E=E_\vartheta-2\beta_B\partial_{\beta_B}
 -2\beta_m\partial_{\beta_m}-\beta_a\partial_{\beta_a}
 -\beta_c\partial_{\beta_c}.
\end{equation}
Holding the original field fixed gives $d/dh=-\mathfrak E/(2h)$.
Analyticity of the principal coefficients and their jet comparison
imply $C/A=2/5+O(|\beta|+r_s)$, with its first two
$\mathfrak E$ derivatives $O(|\beta|+r_s)$. Thus, after fixed
restrictions,
\begin{equation}\label{rw:source:q-rates}
 A>A_*/2,\quad C>C_*/2,\qquad
 q=hC/A,\qquad q_h>1/5,\qquad
 q_{hh}=O((|\beta|+r_s)/h).
\end{equation}
For instance, if $R=C/A$ then
$q_h=R-\mathfrak ER/2$ and
$q_{hh}=(\mathfrak E^2R-2\mathfrak ER)/(4h)$.

\subsection{The exact logarithmic forcing}

Let $U_{\pm,0}=U_\pm|_{\widehat\ell=0}$. Hadamard division of the
signed core equations gives the exact secants
\begin{equation}\label{rw:source:secants}
 \frac{U_+}{U_{+,0}}=1+\widehat\ell\widetilde I_+,
 \qquad
 \frac{U_-}{U_{-,0}}=1-\widehat\ell\widetilde I_-,
 \qquad 1/3<\widetilde I_\pm<3.
\end{equation}
Here $\widetilde I_\pm=I_\pm/U_{\pm,0}
+O(e^{-h/3})$ in the selected independent jets. The actual
fixed-field derivative, before suppressing any independent variable,
is
\begin{equation}\label{rw:source:physical-D}
 D_L:=\frac d{dL}
 =\widehat\ell\partial_{\widehat\ell}-E_s
   -\frac{1-s}{2h}\left(E_\vartheta
       -2\beta_B\partial_{\beta_B}-2\beta_m\partial_{\beta_m}
       -\beta_a\partial_{\beta_a}-\beta_c\partial_{\beta_c}\right).
\end{equation}
In particular $D_L\widehat\ell=\widehat\ell$, $D_Lh=1-s$, $D_L^2h=s$,
and
$D_L\widetilde I_\pm=O(h^{-1})$,
$D_L^2\widetilde I_\pm=O(h^{-2})$. Derivatives of
$U_{\pm,0}-1$ remain exponentially small.

The normalized cores can approach zero, so the tail entries require
the full logarithms
\[
 \zeta_+=\log(1+\widehat\ell\widetilde I_+),\qquad
 \zeta_-=-\log(1-\widehat\ell\widetilde I_-).
\]
They are defined on \eqref{rw:source:positive-core}; for $h$ large,
$|\zeta_\pm|/h<1/1024$. Write
$a_\pm=\widehat\ell\widetilde I_\pm$,
$u_\pm=D_L\widetilde I_\pm/\widetilde I_\pm$ and
$w_\pm=D_L^2\widetilde I_\pm/\widetilde I_\pm$.
Exact differentiation gives
\begin{equation}\label{rw:source:log-jets}
\begin{aligned}
 D_L\zeta_+&=\frac{a_+(1+u_+)}{1+a_+},
 &D_L^2\zeta_+&=
 \frac{a_+[1+2u_++w_++a_+(w_+-u_+^2)]}{(1+a_+)^2},\\
 D_L\zeta_-&=\frac{a_-(1+u_-)}{1-a_-},
 &D_L^2\zeta_-&=
 \frac{a_-[1+2u_-+w_-+a_-(u_-^2-w_-)]}{(1-a_-)^2}.
\end{aligned}
\end{equation}
For $0<\widehat\ell\le3$, their values and first two derivatives are
positive. The smallness of $u_\pm,w_\pm$, $|a_\pm|\le9$ and
elementary logarithm inequalities give fixed constants such that
\begin{equation}\label{rw:source:log-comparison}
 |\zeta_\pm|\le C|D_L\zeta_\pm|\le C^2|D_L^2\zeta_\pm|,
 \quad |D_L\zeta_\pm|^2\le C|D_L^2\zeta_\pm|,
 \quad D_L^j\zeta_\pm\ge c\widehat\ell\quad(j=0,1,2).
\end{equation}
The denominators in \eqref{rw:source:log-jets} are retained in
these comparisons, even if a core value is small.

Apply the unforced incoming tail inverse to the two actions in
\eqref{rw:source:shooting}, keeping its two entries independent.
The divided tail estimate and the bounded $\alpha$ derivative of the
upper map yield an exact identity
\begin{equation}\label{rw:source:independent-forcing}
\begin{split}
 &L_-^{-1}\mathcal U_\alpha T_{+,\widehat\ell}(v_+)
                   -L_-^{-1}T_{-,\widehat\ell}(v_-)\\
 &\hspace{12mm}=\mathcal M(v_+)-v_-
       +\widehat\ell\mathcal E(v_+,v_-,\vartheta,s,\beta,\widehat\ell),
 \qquad |\mathcal E|_N\le C_Ne^{-1/(3\vartheta^2)}.
\end{split}
\end{equation}
The quotient is removable at $\widehat\ell=0$ by analytic dependence.
The actual entries satisfy
$v_+=v_{+,0}-\zeta_+/h$ and
$v_-=v_{-,0}+\zeta_-/h$. After their substitution,
\begin{equation}\label{rw:source:forced-error}
 |D_L^j(h\widehat\ell\mathcal E)|\le C\widehat\ell e^{-h/4},
 \qquad j=0,1,2.
\end{equation}
At most two inverse core factors occur: their cost is bounded by
$s^{-4\delta_0}$, and $1/3-4\delta_0>1/4$ absorbs them and all
fixed polynomial factors of $h$.

Multiply the left side of \eqref{rw:source:independent-forcing}
by $h$ and call the resulting exact function $\mathcal D$. Then
\begin{equation}\label{rw:source:matched}
\begin{split}
 \mathcal D&=\tau_0 f+\sigma_0qf-\Phi,\qquad
 f=h^{3/2}A>0,\quad q=hC/A,\\
 \Phi&=\zeta_+S+\zeta_--h\widehat\ell\mathcal E,\qquad
 S=\int_0^1\mathcal M_v(v_{+,0}-t\zeta_+/h)\,dt.
\end{split}
\end{equation}
This exact section comparison retains the full dependence on
$\widehat\ell$. Its zeros and their multiplicities agree with those
of \eqref{rw:source:shooting}.
On every fixed compact energy interval in its domain, $\Phi/\ell$
has a removable analytic extension at $\ell=0$ with its ordinary
energy derivatives: first divide the independent maps by $\widehat\ell$,
then use $\widehat\ell=\ell/s$ and the positive lower bound of $s$ on
that fixed interval. The same statement holds for any selected
coefficient derivatives after reserving their orders.

\begin{proposition}[Source curvature on one interval]
\label{rw:source:curvature}
There are fixed $\eta>0$ and $H_*<\infty$ such that for every
fixed field in the selected direction $\ell\ge0$, the set
\[
 \mathcal I_s=\{h\ge H_*: |\beta_i|\le\eta,
                  \ 0\le\widehat\ell\le3,
                  \ U_\pm\ge s^{2\delta_0}\}
\]
is an interval, possibly empty. On it $q_h>1/5$.
If $\ell>0$, the function $g=\Phi/f$ satisfies
$g_{qq}>0$. If $\ell=0$, then $\Phi=0$ identically and the
closing equation divided by $f$ is $\tau_0+\sigma_0q=0$.
All statements extend across retained finite endpoints as germs of
analytic section maps.
\end{proposition}
\begin{proof}
At a fixed independent entry, the limiting-map comparison gives
\[
 |\mathcal M_v-1|+|\mathcal M_{vv}|+|\mathcal M_{vvv}|
 +h|D_{L,\mathrm{ind}}\mathcal M_v|
 +h|D_{L,\mathrm{ind}}\mathcal M_{vv}|
 +h^2|D_{L,\mathrm{ind}}^2\mathcal M_v|
 \le C(|\beta|+r_s).
\]
The first two physical derivatives of $v_{+,0}$ have the
corresponding $h^{-1}$ and $h^{-2}$ bounds.
Appendix~\ref{rw:app:source-moving-entry} differentiates the integral
for $S$, keeping also its $\mathcal M_{vvv}(D_Lv)^2$ term.
Using \eqref{rw:source:log-comparison} gives
\[
 |D_L^j\{\zeta_+(S-1)\}|\le C\varepsilon_1
                         |D_L^j\zeta_+|,\qquad j=0,1,2,
\]
where $\varepsilon_1$ is made arbitrarily small by first shrinking
$\eta$ and then increasing $H_*$. Combining this with
\eqref{rw:source:forced-error} proves
\begin{equation}\label{rw:source:Phi-rates}
 \Phi,D_L\Phi,D_L^2\Phi>0,\qquad
 \Phi\le C D_L\Phi\le C^2D_L^2\Phi,
 \qquad D_L\Phi\ge c\widehat\ell.
\end{equation}
In addition \eqref{rw:source:q-rates} gives
$D_Lq\ge c_q>0$, $D_L^2q=O(h^{-1})$,
$D_Lf/f=O(h^{-1})$ and $D_L^2f/f=O(h^{-2})$.
In the exact quotient
\[
 g_{qq}=
 \frac{(D_L^2g)(D_Lq)-(D_Lg)(D_L^2q)}{(D_Lq)^3},
\]
the principal numerator is $(D_Lq)D_L^2\Phi/f>0$.
Every other term is bounded by
$C(|D_L\Phi|/h+|\Phi|/h^2)/f$, less than half of that principal
term when $H_*$ is large. This proves strict curvature.

Before the core restrictions, $h,|\beta_i|,\widehat\ell=\ell e^L$
are nondecreasing functions of $L$, so the permitted energies form
an interval. For $\widehat\ell\ge0$ the outgoing core inequality is
automatic. At any zero of $U_--s^{2\delta_0}$, the secant identity
gives exactly
\[
 D_L(U_--s^{2\delta_0})
 =-U_{-,0}(1+u_-)
 +s^{2\delta_0}\left(1+u_-+2\delta_0
                          +D_LU_{-,0}/U_{-,0}\right)<-1/2.
\]
Thus that boundary can only be crossed downwards, and its retained
side is one interval. At $\ell=0$ both restrictions are automatic.
The strict tube margins prove the endpoint assertion.
\end{proof}

\subsection{Positive return and localization in the chosen direction}

To use the curvature bound for counted cycles, we need a return on
the full interval and a core margin for every such cycle. We first
complete the forward stable core, then use the closing condition to
place each cycle inside the retained positive-core domain.

\begin{proposition}[The actual source return]
\label{rw:source:localization}
Choose the fixed constants and coefficient neighborhood smaller if
necessary. In the direction $\ell\ge0$, the shooting construction
has a positive central first return throughout $\mathcal I_s$, before
its closing equation is imposed. Every counted cycle with
$h\ge H_*$ and $\mathscr M(h)\le\eta$ satisfies the stronger bounds
\begin{equation}\label{rw:source:localization-margin}
 U_+,U_-\ge e^{-\delta_0h}>s^{2\delta_0},
\end{equation}
and is one of its shooting zeros. The positive-return assertion is
for this chosen direction; for an originally negative $\ell$ it
represents the reflected forward, hence original backward, return.
\end{proposition}
\begin{proof}
For $0\le\widehat\ell\le3$, $U_+\ge U_{+,0}>1/2$. The outgoing
entry differs from its unforced value by
$-h^{-1}\log(1+\widehat\ell\widetilde I_+)$, of size at most
$3\widehat\ell/h$. Therefore its positive core, tail and upper arc exist
without any incoming-core assumption. Their negative-axis action is
within $9/1024$ of one. The actual incoming tail inverse, whose domain
was constructed with independent input, gives a unique
$v_{\mathrm{out}}\in[7/16,9/16]$ and a forward endpoint
$x=-\vartheta^{-1}$, $z=e^{-hv_{\mathrm{out}}}>0$.

On the remaining stable core, $-\vartheta^{-1}\le x\le0$,
$0\le z\le1/8$, the original graph denominator exceeds $27/32$.
The field itself gives
\[
 z_x\le2s\widehat\ell+2\vartheta b_0z.
\]
At $z=0$ the normal forcing is nonnegative. For $\ell>0$ its sign
prevents a positive trajectory from crossing downwards; for $\ell=0$
uniqueness preserves positivity. Gronwall on this $x$-interval
gives
\[
 0<z(x)\le C e^{-7h/16}+C\sqrt h\,e^{-h}<1/8.
\]
There is no exit through its upper boundary; the positive horizontal
velocity carries the graph to $x=0$ in finite time. Its height is
below the fixed central-section cutoff after increasing $H_*$.
This is the first return: both lower cores have increasing $x$, the
tails keep their respective signs of $x$, and the upper intersection
with $x=0$ has $y>0$. The core heights tend uniformly to zero,
the tails have $|x|\ge3/(4\vartheta)$ and $-1<y\le0$, and the
upper action has a positive lower bound. Their compactifications
therefore lie in the prescribed collar for sufficiently small scale.

Now consider a counted cycle. Lemma~\ref{rw:geometry:cycle-height}
gives $|\ell|<3s$, so all its independent source parameters belong
to the boxes above. The outgoing construction follows its actual
orbit by uniqueness. At a cycle the returned central height is $s$;
backward uniqueness on the incoming core first proves $U_->0$,
without a quantitative core margin. It also gives
\[
 v_{\mathrm{out}}-v_{-,0}
                  =-h^{-1}\log(U_-/U_{-,0}).
\]
Center division, the bounded upper and tail derivatives, and the
divided forcing estimates imply
\[
 |v_{\mathrm{out}}-v_{-,0}|
 \le C\bigl(|\beta_a+\beta_c|
          +|\beta_c(\beta_B+\beta_m)|+\widehat\ell/h\bigr).
\]
The independently constructed tail inverse makes this estimate
available before a quantitative incoming-core margin. Consequently
\[
 -h^{-1}\log U_-\le C/h+C(2\eta+2\eta^2+3/h).
\]
Choose $\eta$ first to make its contribution less than
$\delta_0/4$, and then $H_*$ large enough for the remaining terms.
The resulting bound is strictly smaller than $\delta_0$.
The outgoing bound follows from $U_+\ge U_{+,0}>1/2$.
This proves \eqref{rw:source:localization-margin} and the claimed
identification with shooting zeros. All constants were selected in
reflection-symmetric boxes; reflection exchanges the two independent
cores and preserves this localization for the original negative sign.
Only the local inverse-return identity near a periodic orbit is needed
later to transfer its multiplicity.
\end{proof}
